
\documentclass{article}
\usepackage{colortbl}
\usepackage{makecell}
\usepackage{multirow}
\usepackage{supertabular}

\begin{document}

\newcounter{utterance}

\centering \large Interaction Transcript for game `codenames', experiment `risk\_low', episode 3 with gemma{-}4{-}26b{-}a4b{-}it.
\vspace{24pt}

{ \footnotesize  \setcounter{utterance}{1}
\setlength{\tabcolsep}{0pt}
\begin{supertabular}{c@{$\;$}|p{.15\linewidth}@{}p{.15\linewidth}p{.15\linewidth}p{.15\linewidth}p{.15\linewidth}p{.15\linewidth}}
    \# & $\;$A & \multicolumn{4}{c}{Game Master} & $\;\:$B\\
    \hline

    \theutterance \stepcounter{utterance}  
    & & & \multicolumn{2}{p{0.3\linewidth}}{
        \cellcolor[rgb]{0.9,0.9,0.9}{
            \makecell[{{p{\linewidth}}}]{
                \texttt{\tiny{[GM$|$GM]}}
                {'hidden': {'team': ['har', 'europa', 'şosetă', 'legătură', 'grevă', 'stâncă', 'soldat', 'inghetata', 'piscină'], 'innocent': ['avion', 'limuzină', 'personal', 'portbagaj', 'teren', 'scaun', 'bici', 'şoim'], 'opponent': ['cupru', 'celulă', 'om de ştiinţă', 'maşină de spălat', 'curea', 'laborator', 'linie', 'pol'], 'assassin': []}, 'revealed': {'team': {'team': [], 'innocent': [], 'opponent': [], 'assassin': []}, 'opponent': {'team': [], 'innocent': [], 'opponent': [], 'assassin': []}}}
            }
        }
    }
    & & \\ \\

    \theutterance \stepcounter{utterance}  
    & & \multicolumn{4}{p{0.6\linewidth}}{
        \cellcolor[rgb]{0.9,0.9,0.9}{
            \makecell[{{p{\linewidth}}}]{
                \texttt{\tiny{[P1$\langle$GM]}}
                \texttt{Furnizați un indiciu dintr{-}un singur cuvânt care se referă la unele dintre cuvintele echipei dvs. din lista de mai jos. Puteți alege la câte și la ce cuvinte se referă indiciul dumneavoastră.} \\
\texttt{Scopul este ca partenerul tău să ghicească toate cuvintele echipei tale pe baza indiciilor pe care i le dai, înainte ca cealaltă echipă să le descopere pe toate. Dacă partenerul dvs. ghicește un cuvânt "adversar" sau un cuvânt "distractor", rândul său de ghicit se încheie, iar echipa adversă are ocazia să facă un tur. Dacă partenerul tău ghicește un cuvânt "ucigaș", jocul se încheie și echipa ta pierde imediat.} \\
\texttt{Cu cât indiciul tău se referă la mai multe cuvinte, cu atât câștigi mai repede jocul, dar cu atât asocierile pot fi mai greu de ghicit pentru partenerul tău, așa că alege{-}ți cu înțelepciune cuvântul indiciu. Cuvântul indiciu trebuie să fie legat semantic de cuvintele țintă, nu poate fi unul dintre cuvintele din liste sau să conțină părți ale acestora.} \\
\\ 
\texttt{Furnizați întotdeauna indiciul cu un singur cuvânt și lista de cuvinte țintă, separate prin virgulă, în următorul format și faceți răspunsurile cât mai scurte posibil; nu includeți niciodată alt text decât este necesar în acest formular:} \\
\\ 
\texttt{INDICIU: <CUVÂNT>} \\
\texttt{ȚINTE: <ȚINTE>} \\
\\ 
\texttt{Cuvintele echipei tale sunt: har, europa, şosetă, legătură, grevă, stâncă, soldat, inghetata, piscină.} \\
\texttt{Cuvintele echipei adversarului dvs. sunt: cupru, celulă, om de ştiinţă, maşină de spălat, curea, laborator, linie, pol.} \\
\texttt{Cuvintele de distragere a atenției sunt: avion, limuzină, personal, portbagaj, teren, scaun, bici, şoim.} \\
\texttt{Cuvintele de ucidere sunt: .} \\
            }
        }
    }
    & \\ \\

    \theutterance \stepcounter{utterance}  
    & \multicolumn{4}{p{0.6\linewidth}}{
        \cellcolor[rgb]{0.9,0.9,0.9}{
            \makecell[{{p{\linewidth}}}]{
                \texttt{\tiny{[P1$\rangle$GM]}}
                \texttt{INDICIU: CONTINENT} \\
\texttt{ȚINTE: europa, legătură} \\
            }
        }
    }
    & & \\ \\

    \theutterance \stepcounter{utterance}  
    & & & \multicolumn{2}{p{0.3\linewidth}}{
        \cellcolor[rgb]{0.9,0.9,0.9}{
            \makecell[{{p{\linewidth}}}]{
                \texttt{\tiny{[GM$|$GM]}}
                \texttt{continent} \\
            }
        }
    }
    & & \\ \\

    \theutterance \stepcounter{utterance}  
    & & & \multicolumn{2}{p{0.3\linewidth}}{
        \cellcolor[rgb]{0.9,0.9,0.9}{
            \makecell[{{p{\linewidth}}}]{
                \texttt{\tiny{[GM$|$GM]}}
                ['europa', 'legătură']
            }
        }
    }
    & & \\ \\

    \theutterance \stepcounter{utterance}  
    & & & \multicolumn{2}{p{0.3\linewidth}}{
        \cellcolor[rgb]{0.9,0.9,0.9}{
            \makecell[{{p{\linewidth}}}]{
                \texttt{\tiny{[GM$|$GM]}}
                {'word': 'europa', 'assignment': 'team'}
            }
        }
    }
    & & \\ \\

    \theutterance \stepcounter{utterance}  
    & & & \multicolumn{2}{p{0.3\linewidth}}{
        \cellcolor[rgb]{0.9,0.9,0.9}{
            \makecell[{{p{\linewidth}}}]{
                \texttt{\tiny{[GM$|$GM]}}
                {'word': 'legătură', 'assignment': 'team'}
            }
        }
    }
    & & \\ \\

    \theutterance \stepcounter{utterance}  
    & & \multicolumn{4}{p{0.6\linewidth}}{
        \cellcolor[rgb]{0.9,0.9,0.9}{
            \makecell[{{p{\linewidth}}}]{
                \texttt{\tiny{[GM$\rangle$P2]}}
                \texttt{Furnizați o listă separată prin virgulă de până la 2 cuvinte din următoarea listă care se referă cel mai bine sau sunt asociate cel mai strâns cu cuvântul 'continent'. Începeți întotdeauna lista de presupuneri cu GHICIRE: și nu includeți niciun alt text în răspunsul dumneavoastră.} \\
\\ 
\texttt{grevă, europa, maşină de spălat, portbagaj, şosetă, curea, bici, avion, linie, piscină, şoim, stâncă, personal, har, legătură, laborator, scaun, pol, inghetata, teren, celulă, soldat, cupru, om de ştiinţă, limuzină} \\
            }
        }
    }
    & \\ \\

    \theutterance \stepcounter{utterance}  
    & & & \multicolumn{4}{p{0.6\linewidth}}{
        \cellcolor[rgb]{0.9,0.9,0.9}{
            \makecell[{{p{\linewidth}}}]{
                \texttt{\tiny{[GM$\langle$P2]}}
                \texttt{GHICIRE: europa, pol} \\
            }
        }
    }
     \\ \\

    \theutterance \stepcounter{utterance}  
    & & & \multicolumn{2}{p{0.3\linewidth}}{
        \cellcolor[rgb]{0.9,0.9,0.9}{
            \makecell[{{p{\linewidth}}}]{
                \texttt{\tiny{[GM$|$GM]}}
                ['europa', 'pol']
            }
        }
    }
    & & \\ \\

    \theutterance \stepcounter{utterance}  
    & & & \multicolumn{2}{p{0.3\linewidth}}{
        \cellcolor[rgb]{0.9,0.9,0.9}{
            \makecell[{{p{\linewidth}}}]{
                \texttt{\tiny{[GM$|$GM]}}
                {'word': 'europa', 'assignment': 'team'}
            }
        }
    }
    & & \\ \\

    \theutterance \stepcounter{utterance}  
    & & & \multicolumn{2}{p{0.3\linewidth}}{
        \cellcolor[rgb]{0.9,0.9,0.9}{
            \makecell[{{p{\linewidth}}}]{
                \texttt{\tiny{[GM$|$GM]}}
                {'word': 'europa', 'assignment': 'team'}
            }
        }
    }
    & & \\ \\

    \theutterance \stepcounter{utterance}  
    & & & \multicolumn{2}{p{0.3\linewidth}}{
        \cellcolor[rgb]{0.9,0.9,0.9}{
            \makecell[{{p{\linewidth}}}]{
                \texttt{\tiny{[GM$|$GM]}}
                {'word': 'pol', 'assignment': 'opponent'}
            }
        }
    }
    & & \\ \\

    \theutterance \stepcounter{utterance}  
    & & & \multicolumn{2}{p{0.3\linewidth}}{
        \cellcolor[rgb]{0.9,0.9,0.9}{
            \makecell[{{p{\linewidth}}}]{
                \texttt{\tiny{[GM$|$GM]}}
                \texttt{pol} \\
            }
        }
    }
    & & \\ \\

    \theutterance \stepcounter{utterance}  
    & & & \multicolumn{2}{p{0.3\linewidth}}{
        \cellcolor[rgb]{0.9,0.9,0.9}{
            \makecell[{{p{\linewidth}}}]{
                \texttt{\tiny{[GM$|$GM]}}
                {'word': 'laborator', 'assignment': 'opponent'}
            }
        }
    }
    & & \\ \\

    \theutterance \stepcounter{utterance}  
    & & & \multicolumn{2}{p{0.3\linewidth}}{
        \cellcolor[rgb]{0.9,0.9,0.9}{
            \makecell[{{p{\linewidth}}}]{
                \texttt{\tiny{[GM$|$GM]}}
                {'hidden': {'team': ['har', 'şosetă', 'legătură', 'grevă', 'stâncă', 'soldat', 'inghetata', 'piscină'], 'innocent': ['avion', 'limuzină', 'personal', 'portbagaj', 'teren', 'scaun', 'bici', 'şoim'], 'opponent': ['cupru', 'celulă', 'om de ştiinţă', 'maşină de spălat', 'curea', 'linie'], 'assassin': []}, 'revealed': {'team': {'team': ['europa'], 'innocent': [], 'opponent': ['pol'], 'assassin': []}, 'opponent': {'team': [], 'innocent': [], 'opponent': ['laborator'], 'assassin': []}}}
            }
        }
    }
    & & \\ \\

    \theutterance \stepcounter{utterance}  
    & & \multicolumn{4}{p{0.6\linewidth}}{
        \cellcolor[rgb]{0.9,0.9,0.9}{
            \makecell[{{p{\linewidth}}}]{
                \texttt{\tiny{[P1$\langle$GM]}}
                \texttt{Cuvântul europa a fost ghicit corect. Cuvântul pol a fost ghicit, dar este un cuvânt adversar. Rândul colegului tău de echipă s{-}a încheiat acolo.} \\
\texttt{Acum furnizați un alt indiciu referitor la unele dintre cuvintele echipei dvs. rămase și o listă a cuvintelor țintă aferente. Nu uitați să începeți indiciul cu 'INDICIU: ', puneți un rând nou și începeți lista de cuvinte țintă separate prin virgulă cu 'ȚINTE: '. Observație: unele cuvinte au fost eliminate din liste față de solicitările anterioare.} \\
\\ 
\texttt{Cuvintele echipei dvs. rămase sunt: har, şosetă, legătură, grevă, stâncă, soldat, inghetata, piscină.} \\
\texttt{Cuvintele rămase pentru adversarul dvs. sunt: cupru, celulă, om de ştiinţă, maşină de spălat, curea, laborator, linie.} \\
\texttt{Cuvintele de distragere a atenției rămase sunt: avion, limuzină, personal, portbagaj, teren, scaun, bici, şoim.} \\
\texttt{Cuvintele de ucidere rămase sunt: .} \\
            }
        }
    }
    & \\ \\

    \theutterance \stepcounter{utterance}  
    & \multicolumn{4}{p{0.6\linewidth}}{
        \cellcolor[rgb]{0.9,0.9,0.9}{
            \makecell[{{p{\linewidth}}}]{
                \texttt{\tiny{[P1$\rangle$GM]}}
                \texttt{INDICIU: APĂ} \\
\texttt{ȚINTE: piscină, inghetata} \\
            }
        }
    }
    & & \\ \\

    \theutterance \stepcounter{utterance}  
    & & & \multicolumn{2}{p{0.3\linewidth}}{
        \cellcolor[rgb]{0.9,0.9,0.9}{
            \makecell[{{p{\linewidth}}}]{
                \texttt{\tiny{[GM$|$GM]}}
                \texttt{apă} \\
            }
        }
    }
    & & \\ \\

    \theutterance \stepcounter{utterance}  
    & & & \multicolumn{2}{p{0.3\linewidth}}{
        \cellcolor[rgb]{0.9,0.9,0.9}{
            \makecell[{{p{\linewidth}}}]{
                \texttt{\tiny{[GM$|$GM]}}
                ['piscină', 'inghetata']
            }
        }
    }
    & & \\ \\

    \theutterance \stepcounter{utterance}  
    & & & \multicolumn{2}{p{0.3\linewidth}}{
        \cellcolor[rgb]{0.9,0.9,0.9}{
            \makecell[{{p{\linewidth}}}]{
                \texttt{\tiny{[GM$|$GM]}}
                {'word': 'piscină', 'assignment': 'team'}
            }
        }
    }
    & & \\ \\

    \theutterance \stepcounter{utterance}  
    & & & \multicolumn{2}{p{0.3\linewidth}}{
        \cellcolor[rgb]{0.9,0.9,0.9}{
            \makecell[{{p{\linewidth}}}]{
                \texttt{\tiny{[GM$|$GM]}}
                {'word': 'inghetata', 'assignment': 'team'}
            }
        }
    }
    & & \\ \\

    \theutterance \stepcounter{utterance}  
    & & \multicolumn{4}{p{0.6\linewidth}}{
        \cellcolor[rgb]{0.9,0.9,0.9}{
            \makecell[{{p{\linewidth}}}]{
                \texttt{\tiny{[GM$\rangle$P2]}}
                \texttt{Cuvântul europa a fost ghicit corect. Cuvântul pol a fost ghicit, dar este un cuvânt adversar. Rândul tău s{-}a încheiat acolo.} \\
\texttt{Acum furnizați o altă listă separată prin virgulă de cel puțin 1 și până la 2 cuvinte din următoarea listă de cuvinte care se referă cel mai bine sau sunt cel mai strâns asociate cu cuvântul 'apă'. Nu uitați să vă începeți răspunsul cu GHICIRE:. Notă: unele cuvinte au fost eliminate din listă față de solicitările anterioare.} \\
\\ 
\texttt{grevă, maşină de spălat, portbagaj, şosetă, curea, bici, avion, linie, piscină, şoim, stâncă, personal, har, legătură, scaun, inghetata, teren, celulă, soldat, cupru, om de ştiinţă, limuzină} \\
            }
        }
    }
    & \\ \\

    \theutterance \stepcounter{utterance}  
    & & & \multicolumn{4}{p{0.6\linewidth}}{
        \cellcolor[rgb]{0.9,0.9,0.9}{
            \makecell[{{p{\linewidth}}}]{
                \texttt{\tiny{[GM$\langle$P2]}}
                \texttt{GHICIRE: piscină, inghetata} \\
            }
        }
    }
     \\ \\

    \theutterance \stepcounter{utterance}  
    & & & \multicolumn{2}{p{0.3\linewidth}}{
        \cellcolor[rgb]{0.9,0.9,0.9}{
            \makecell[{{p{\linewidth}}}]{
                \texttt{\tiny{[GM$|$GM]}}
                ['piscină', 'inghetata']
            }
        }
    }
    & & \\ \\

    \theutterance \stepcounter{utterance}  
    & & & \multicolumn{2}{p{0.3\linewidth}}{
        \cellcolor[rgb]{0.9,0.9,0.9}{
            \makecell[{{p{\linewidth}}}]{
                \texttt{\tiny{[GM$|$GM]}}
                {'word': 'piscină', 'assignment': 'team'}
            }
        }
    }
    & & \\ \\

    \theutterance \stepcounter{utterance}  
    & & & \multicolumn{2}{p{0.3\linewidth}}{
        \cellcolor[rgb]{0.9,0.9,0.9}{
            \makecell[{{p{\linewidth}}}]{
                \texttt{\tiny{[GM$|$GM]}}
                {'word': 'piscină', 'assignment': 'team'}
            }
        }
    }
    & & \\ \\

    \theutterance \stepcounter{utterance}  
    & & & \multicolumn{2}{p{0.3\linewidth}}{
        \cellcolor[rgb]{0.9,0.9,0.9}{
            \makecell[{{p{\linewidth}}}]{
                \texttt{\tiny{[GM$|$GM]}}
                {'word': 'inghetata', 'assignment': 'team'}
            }
        }
    }
    & & \\ \\

    \theutterance \stepcounter{utterance}  
    & & & \multicolumn{2}{p{0.3\linewidth}}{
        \cellcolor[rgb]{0.9,0.9,0.9}{
            \makecell[{{p{\linewidth}}}]{
                \texttt{\tiny{[GM$|$GM]}}
                {'word': 'inghetata', 'assignment': 'team'}
            }
        }
    }
    & & \\ \\

    \theutterance \stepcounter{utterance}  
    & & & \multicolumn{2}{p{0.3\linewidth}}{
        \cellcolor[rgb]{0.9,0.9,0.9}{
            \makecell[{{p{\linewidth}}}]{
                \texttt{\tiny{[GM$|$GM]}}
                {'word': 'maşină de spălat', 'assignment': 'opponent'}
            }
        }
    }
    & & \\ \\

    \theutterance \stepcounter{utterance}  
    & & & \multicolumn{2}{p{0.3\linewidth}}{
        \cellcolor[rgb]{0.9,0.9,0.9}{
            \makecell[{{p{\linewidth}}}]{
                \texttt{\tiny{[GM$|$GM]}}
                {'hidden': {'team': ['har', 'şosetă', 'legătură', 'grevă', 'stâncă', 'soldat'], 'innocent': ['avion', 'limuzină', 'personal', 'portbagaj', 'teren', 'scaun', 'bici', 'şoim'], 'opponent': ['cupru', 'celulă', 'om de ştiinţă', 'curea', 'linie'], 'assassin': []}, 'revealed': {'team': {'team': ['europa', 'piscină', 'inghetata'], 'innocent': [], 'opponent': ['pol'], 'assassin': []}, 'opponent': {'team': [], 'innocent': [], 'opponent': ['laborator', 'maşină de spălat'], 'assassin': []}}}
            }
        }
    }
    & & \\ \\

    \theutterance \stepcounter{utterance}  
    & & \multicolumn{4}{p{0.6\linewidth}}{
        \cellcolor[rgb]{0.9,0.9,0.9}{
            \makecell[{{p{\linewidth}}}]{
                \texttt{\tiny{[P1$\langle$GM]}}
                \texttt{Cuvintele piscină, inghetata au fost ghicite corect. Rândul colegului tău de echipă s{-}a încheiat acolo.} \\
\texttt{Acum furnizați un alt indiciu referitor la unele dintre cuvintele echipei dvs. rămase și o listă a cuvintelor țintă aferente. Nu uitați să începeți indiciul cu 'INDICIU: ', puneți un rând nou și începeți lista de cuvinte țintă separate prin virgulă cu 'ȚINTE: '. Observație: unele cuvinte au fost eliminate din liste față de solicitările anterioare.} \\
\\ 
\texttt{Cuvintele echipei dvs. rămase sunt: har, şosetă, legătură, grevă, stâncă, soldat.} \\
\texttt{Cuvintele rămase pentru adversarul dvs. sunt: cupru, celulă, om de ştiinţă, maşină de spălat, curea, linie.} \\
\texttt{Cuvintele de distragere a atenției rămase sunt: avion, limuzină, personal, portbagaj, teren, scaun, bici, şoim.} \\
\texttt{Cuvintele de ucidere rămase sunt: .} \\
            }
        }
    }
    & \\ \\

    \theutterance \stepcounter{utterance}  
    & \multicolumn{4}{p{0.6\linewidth}}{
        \cellcolor[rgb]{0.9,0.9,0.9}{
            \makecell[{{p{\linewidth}}}]{
                \texttt{\tiny{[P1$\rangle$GM]}}
                \texttt{INDICIU: CONFLICT} \\
\texttt{ȚINTE: grevă, soldat} \\
            }
        }
    }
    & & \\ \\

    \theutterance \stepcounter{utterance}  
    & & & \multicolumn{2}{p{0.3\linewidth}}{
        \cellcolor[rgb]{0.9,0.9,0.9}{
            \makecell[{{p{\linewidth}}}]{
                \texttt{\tiny{[GM$|$GM]}}
                \texttt{conflict} \\
            }
        }
    }
    & & \\ \\

    \theutterance \stepcounter{utterance}  
    & & & \multicolumn{2}{p{0.3\linewidth}}{
        \cellcolor[rgb]{0.9,0.9,0.9}{
            \makecell[{{p{\linewidth}}}]{
                \texttt{\tiny{[GM$|$GM]}}
                ['grevă', 'soldat']
            }
        }
    }
    & & \\ \\

    \theutterance \stepcounter{utterance}  
    & & & \multicolumn{2}{p{0.3\linewidth}}{
        \cellcolor[rgb]{0.9,0.9,0.9}{
            \makecell[{{p{\linewidth}}}]{
                \texttt{\tiny{[GM$|$GM]}}
                {'word': 'grevă', 'assignment': 'team'}
            }
        }
    }
    & & \\ \\

    \theutterance \stepcounter{utterance}  
    & & & \multicolumn{2}{p{0.3\linewidth}}{
        \cellcolor[rgb]{0.9,0.9,0.9}{
            \makecell[{{p{\linewidth}}}]{
                \texttt{\tiny{[GM$|$GM]}}
                {'word': 'soldat', 'assignment': 'team'}
            }
        }
    }
    & & \\ \\

    \theutterance \stepcounter{utterance}  
    & & \multicolumn{4}{p{0.6\linewidth}}{
        \cellcolor[rgb]{0.9,0.9,0.9}{
            \makecell[{{p{\linewidth}}}]{
                \texttt{\tiny{[GM$\rangle$P2]}}
                \texttt{Cuvintele piscină, inghetata au fost ghicite corect. Rândul tău s{-}a încheiat acolo.} \\
\texttt{Acum furnizați o altă listă separată prin virgulă de cel puțin 1 și până la 2 cuvinte din următoarea listă de cuvinte care se referă cel mai bine sau sunt cel mai strâns asociate cu cuvântul 'conflict'. Nu uitați să vă începeți răspunsul cu GHICIRE:. Notă: unele cuvinte au fost eliminate din listă față de solicitările anterioare.} \\
\\ 
\texttt{grevă, portbagaj, şosetă, curea, bici, avion, linie, şoim, stâncă, personal, har, legătură, scaun, teren, celulă, soldat, cupru, om de ştiinţă, limuzină} \\
            }
        }
    }
    & \\ \\

    \theutterance \stepcounter{utterance}  
    & & & \multicolumn{4}{p{0.6\linewidth}}{
        \cellcolor[rgb]{0.9,0.9,0.9}{
            \makecell[{{p{\linewidth}}}]{
                \texttt{\tiny{[GM$\langle$P2]}}
                \texttt{GHICIRE: grevă, soldat} \\
            }
        }
    }
     \\ \\

    \theutterance \stepcounter{utterance}  
    & & & \multicolumn{2}{p{0.3\linewidth}}{
        \cellcolor[rgb]{0.9,0.9,0.9}{
            \makecell[{{p{\linewidth}}}]{
                \texttt{\tiny{[GM$|$GM]}}
                ['grevă', 'soldat']
            }
        }
    }
    & & \\ \\

    \theutterance \stepcounter{utterance}  
    & & & \multicolumn{2}{p{0.3\linewidth}}{
        \cellcolor[rgb]{0.9,0.9,0.9}{
            \makecell[{{p{\linewidth}}}]{
                \texttt{\tiny{[GM$|$GM]}}
                {'word': 'grevă', 'assignment': 'team'}
            }
        }
    }
    & & \\ \\

    \theutterance \stepcounter{utterance}  
    & & & \multicolumn{2}{p{0.3\linewidth}}{
        \cellcolor[rgb]{0.9,0.9,0.9}{
            \makecell[{{p{\linewidth}}}]{
                \texttt{\tiny{[GM$|$GM]}}
                {'word': 'grevă', 'assignment': 'team'}
            }
        }
    }
    & & \\ \\

    \theutterance \stepcounter{utterance}  
    & & & \multicolumn{2}{p{0.3\linewidth}}{
        \cellcolor[rgb]{0.9,0.9,0.9}{
            \makecell[{{p{\linewidth}}}]{
                \texttt{\tiny{[GM$|$GM]}}
                {'word': 'soldat', 'assignment': 'team'}
            }
        }
    }
    & & \\ \\

    \theutterance \stepcounter{utterance}  
    & & & \multicolumn{2}{p{0.3\linewidth}}{
        \cellcolor[rgb]{0.9,0.9,0.9}{
            \makecell[{{p{\linewidth}}}]{
                \texttt{\tiny{[GM$|$GM]}}
                {'word': 'soldat', 'assignment': 'team'}
            }
        }
    }
    & & \\ \\

    \theutterance \stepcounter{utterance}  
    & & & \multicolumn{2}{p{0.3\linewidth}}{
        \cellcolor[rgb]{0.9,0.9,0.9}{
            \makecell[{{p{\linewidth}}}]{
                \texttt{\tiny{[GM$|$GM]}}
                {'word': 'celulă', 'assignment': 'opponent'}
            }
        }
    }
    & & \\ \\

    \theutterance \stepcounter{utterance}  
    & & & \multicolumn{2}{p{0.3\linewidth}}{
        \cellcolor[rgb]{0.9,0.9,0.9}{
            \makecell[{{p{\linewidth}}}]{
                \texttt{\tiny{[GM$|$GM]}}
                {'hidden': {'team': ['har', 'şosetă', 'legătură', 'stâncă'], 'innocent': ['avion', 'limuzină', 'personal', 'portbagaj', 'teren', 'scaun', 'bici', 'şoim'], 'opponent': ['cupru', 'om de ştiinţă', 'curea', 'linie'], 'assassin': []}, 'revealed': {'team': {'team': ['europa', 'piscină', 'inghetata', 'grevă', 'soldat'], 'innocent': [], 'opponent': ['pol'], 'assassin': []}, 'opponent': {'team': [], 'innocent': [], 'opponent': ['laborator', 'maşină de spălat', 'celulă'], 'assassin': []}}}
            }
        }
    }
    & & \\ \\

    \theutterance \stepcounter{utterance}  
    & & \multicolumn{4}{p{0.6\linewidth}}{
        \cellcolor[rgb]{0.9,0.9,0.9}{
            \makecell[{{p{\linewidth}}}]{
                \texttt{\tiny{[P1$\langle$GM]}}
                \texttt{Cuvintele grevă, soldat au fost ghicite corect. Rândul colegului tău de echipă s{-}a încheiat acolo.} \\
\texttt{Acum furnizați un alt indiciu referitor la unele dintre cuvintele echipei dvs. rămase și o listă a cuvintelor țintă aferente. Nu uitați să începeți indiciul cu 'INDICIU: ', puneți un rând nou și începeți lista de cuvinte țintă separate prin virgulă cu 'ȚINTE: '. Observație: unele cuvinte au fost eliminate din liste față de solicitările anterioare.} \\
\\ 
\texttt{Cuvintele echipei dvs. rămase sunt: har, şosetă, legătură, stâncă.} \\
\texttt{Cuvintele rămase pentru adversarul dvs. sunt: cupru, celulă, om de ştiinţă, curea, linie.} \\
\texttt{Cuvintele de distragere a atenției rămase sunt: avion, limuzină, personal, portbagaj, teren, scaun, bici, şoim.} \\
\texttt{Cuvintele de ucidere rămase sunt: .} \\
            }
        }
    }
    & \\ \\

    \theutterance \stepcounter{utterance}  
    & \multicolumn{4}{p{0.6\linewidth}}{
        \cellcolor[rgb]{0.9,0.9,0.9}{
            \makecell[{{p{\linewidth}}}]{
                \texttt{\tiny{[P1$\rangle$GM]}}
                \texttt{INDICIU: PIATRĂ} \\
\texttt{ȚINTE: stâncă, har} \\
            }
        }
    }
    & & \\ \\

    \theutterance \stepcounter{utterance}  
    & & & \multicolumn{2}{p{0.3\linewidth}}{
        \cellcolor[rgb]{0.9,0.9,0.9}{
            \makecell[{{p{\linewidth}}}]{
                \texttt{\tiny{[GM$|$GM]}}
                \texttt{piatră} \\
            }
        }
    }
    & & \\ \\

    \theutterance \stepcounter{utterance}  
    & & & \multicolumn{2}{p{0.3\linewidth}}{
        \cellcolor[rgb]{0.9,0.9,0.9}{
            \makecell[{{p{\linewidth}}}]{
                \texttt{\tiny{[GM$|$GM]}}
                ['stâncă', 'har']
            }
        }
    }
    & & \\ \\

    \theutterance \stepcounter{utterance}  
    & & & \multicolumn{2}{p{0.3\linewidth}}{
        \cellcolor[rgb]{0.9,0.9,0.9}{
            \makecell[{{p{\linewidth}}}]{
                \texttt{\tiny{[GM$|$GM]}}
                {'word': 'stâncă', 'assignment': 'team'}
            }
        }
    }
    & & \\ \\

    \theutterance \stepcounter{utterance}  
    & & & \multicolumn{2}{p{0.3\linewidth}}{
        \cellcolor[rgb]{0.9,0.9,0.9}{
            \makecell[{{p{\linewidth}}}]{
                \texttt{\tiny{[GM$|$GM]}}
                {'word': 'har', 'assignment': 'team'}
            }
        }
    }
    & & \\ \\

    \theutterance \stepcounter{utterance}  
    & & \multicolumn{4}{p{0.6\linewidth}}{
        \cellcolor[rgb]{0.9,0.9,0.9}{
            \makecell[{{p{\linewidth}}}]{
                \texttt{\tiny{[GM$\rangle$P2]}}
                \texttt{Cuvintele grevă, soldat au fost ghicite corect. Rândul tău s{-}a încheiat acolo.} \\
\texttt{Acum furnizați o altă listă separată prin virgulă de cel puțin 1 și până la 2 cuvinte din următoarea listă de cuvinte care se referă cel mai bine sau sunt cel mai strâns asociate cu cuvântul 'piatră'. Nu uitați să vă începeți răspunsul cu GHICIRE:. Notă: unele cuvinte au fost eliminate din listă față de solicitările anterioare.} \\
\\ 
\texttt{portbagaj, şosetă, curea, bici, avion, linie, şoim, stâncă, personal, har, legătură, scaun, teren, cupru, om de ştiinţă, limuzină} \\
            }
        }
    }
    & \\ \\

    \theutterance \stepcounter{utterance}  
    & & & \multicolumn{4}{p{0.6\linewidth}}{
        \cellcolor[rgb]{0.9,0.9,0.9}{
            \makecell[{{p{\linewidth}}}]{
                \texttt{\tiny{[GM$\langle$P2]}}
                \texttt{GHICIRE: stâncă, teren} \\
            }
        }
    }
     \\ \\

    \theutterance \stepcounter{utterance}  
    & & & \multicolumn{2}{p{0.3\linewidth}}{
        \cellcolor[rgb]{0.9,0.9,0.9}{
            \makecell[{{p{\linewidth}}}]{
                \texttt{\tiny{[GM$|$GM]}}
                ['stâncă', 'teren']
            }
        }
    }
    & & \\ \\

    \theutterance \stepcounter{utterance}  
    & & & \multicolumn{2}{p{0.3\linewidth}}{
        \cellcolor[rgb]{0.9,0.9,0.9}{
            \makecell[{{p{\linewidth}}}]{
                \texttt{\tiny{[GM$|$GM]}}
                {'word': 'stâncă', 'assignment': 'team'}
            }
        }
    }
    & & \\ \\

    \theutterance \stepcounter{utterance}  
    & & & \multicolumn{2}{p{0.3\linewidth}}{
        \cellcolor[rgb]{0.9,0.9,0.9}{
            \makecell[{{p{\linewidth}}}]{
                \texttt{\tiny{[GM$|$GM]}}
                {'word': 'stâncă', 'assignment': 'team'}
            }
        }
    }
    & & \\ \\

    \theutterance \stepcounter{utterance}  
    & & & \multicolumn{2}{p{0.3\linewidth}}{
        \cellcolor[rgb]{0.9,0.9,0.9}{
            \makecell[{{p{\linewidth}}}]{
                \texttt{\tiny{[GM$|$GM]}}
                {'word': 'teren', 'assignment': 'innocent'}
            }
        }
    }
    & & \\ \\

    \theutterance \stepcounter{utterance}  
    & & & \multicolumn{2}{p{0.3\linewidth}}{
        \cellcolor[rgb]{0.9,0.9,0.9}{
            \makecell[{{p{\linewidth}}}]{
                \texttt{\tiny{[GM$|$GM]}}
                \texttt{teren} \\
            }
        }
    }
    & & \\ \\

    \theutterance \stepcounter{utterance}  
    & & & \multicolumn{2}{p{0.3\linewidth}}{
        \cellcolor[rgb]{0.9,0.9,0.9}{
            \makecell[{{p{\linewidth}}}]{
                \texttt{\tiny{[GM$|$GM]}}
                {'word': 'curea', 'assignment': 'opponent'}
            }
        }
    }
    & & \\ \\

    \theutterance \stepcounter{utterance}  
    & & & \multicolumn{2}{p{0.3\linewidth}}{
        \cellcolor[rgb]{0.9,0.9,0.9}{
            \makecell[{{p{\linewidth}}}]{
                \texttt{\tiny{[GM$|$GM]}}
                {'hidden': {'team': ['har', 'şosetă', 'legătură'], 'innocent': ['avion', 'limuzină', 'personal', 'portbagaj', 'scaun', 'bici', 'şoim'], 'opponent': ['cupru', 'om de ştiinţă', 'linie'], 'assassin': []}, 'revealed': {'team': {'team': ['europa', 'piscină', 'inghetata', 'grevă', 'soldat', 'stâncă'], 'innocent': ['teren'], 'opponent': ['pol'], 'assassin': []}, 'opponent': {'team': [], 'innocent': [], 'opponent': ['laborator', 'maşină de spălat', 'celulă', 'curea'], 'assassin': []}}}
            }
        }
    }
    & & \\ \\

    \theutterance \stepcounter{utterance}  
    & & \multicolumn{4}{p{0.6\linewidth}}{
        \cellcolor[rgb]{0.9,0.9,0.9}{
            \makecell[{{p{\linewidth}}}]{
                \texttt{\tiny{[P1$\langle$GM]}}
                \texttt{Cuvântul stâncă a fost ghicit corect. Cuvântul teren a fost ghicit, dar este un cuvânt nevinovat. Rândul colegului tău de echipă s{-}a încheiat acolo.} \\
\texttt{Acum furnizați un alt indiciu referitor la unele dintre cuvintele echipei dvs. rămase și o listă a cuvintelor țintă aferente. Nu uitați să începeți indiciul cu 'INDICIU: ', puneți un rând nou și începeți lista de cuvinte țintă separate prin virgulă cu 'ȚINTE: '. Observație: unele cuvinte au fost eliminate din liste față de solicitările anterioare.} \\
\\ 
\texttt{Cuvintele echipei dvs. rămase sunt: har, şosetă, legătură.} \\
\texttt{Cuvintele rămase pentru adversarul dvs. sunt: cupru, om de ştiinţă, curea, linie.} \\
\texttt{Cuvintele de distragere a atenției rămase sunt: avion, limuzină, personal, portbagaj, scaun, bici, şoim.} \\
\texttt{Cuvintele de ucidere rămase sunt: .} \\
            }
        }
    }
    & \\ \\

    \theutterance \stepcounter{utterance}  
    & \multicolumn{4}{p{0.6\linewidth}}{
        \cellcolor[rgb]{0.9,0.9,0.9}{
            \makecell[{{p{\linewidth}}}]{
                \texttt{\tiny{[P1$\rangle$GM]}}
                \texttt{INDICIU: PANTALONI} \\
\texttt{ȚINTE: şosetă, legătură} \\
            }
        }
    }
    & & \\ \\

    \theutterance \stepcounter{utterance}  
    & & & \multicolumn{2}{p{0.3\linewidth}}{
        \cellcolor[rgb]{0.9,0.9,0.9}{
            \makecell[{{p{\linewidth}}}]{
                \texttt{\tiny{[GM$|$GM]}}
                \texttt{pantaloni} \\
            }
        }
    }
    & & \\ \\

    \theutterance \stepcounter{utterance}  
    & & & \multicolumn{2}{p{0.3\linewidth}}{
        \cellcolor[rgb]{0.9,0.9,0.9}{
            \makecell[{{p{\linewidth}}}]{
                \texttt{\tiny{[GM$|$GM]}}
                ['şosetă', 'legătură']
            }
        }
    }
    & & \\ \\

    \theutterance \stepcounter{utterance}  
    & & & \multicolumn{2}{p{0.3\linewidth}}{
        \cellcolor[rgb]{0.9,0.9,0.9}{
            \makecell[{{p{\linewidth}}}]{
                \texttt{\tiny{[GM$|$GM]}}
                {'word': 'şosetă', 'assignment': 'team'}
            }
        }
    }
    & & \\ \\

    \theutterance \stepcounter{utterance}  
    & & & \multicolumn{2}{p{0.3\linewidth}}{
        \cellcolor[rgb]{0.9,0.9,0.9}{
            \makecell[{{p{\linewidth}}}]{
                \texttt{\tiny{[GM$|$GM]}}
                {'word': 'legătură', 'assignment': 'team'}
            }
        }
    }
    & & \\ \\

    \theutterance \stepcounter{utterance}  
    & & \multicolumn{4}{p{0.6\linewidth}}{
        \cellcolor[rgb]{0.9,0.9,0.9}{
            \makecell[{{p{\linewidth}}}]{
                \texttt{\tiny{[GM$\rangle$P2]}}
                \texttt{Cuvântul stâncă a fost ghicit corect. Cuvântul teren a fost ghicit, dar este un cuvânt nevinovat. Rândul tău s{-}a încheiat acolo.} \\
\texttt{Acum furnizați o altă listă separată prin virgulă de cel puțin 1 și până la 2 cuvinte din următoarea listă de cuvinte care se referă cel mai bine sau sunt cel mai strâns asociate cu cuvântul 'pantaloni'. Nu uitați să vă începeți răspunsul cu GHICIRE:. Notă: unele cuvinte au fost eliminate din listă față de solicitările anterioare.} \\
\\ 
\texttt{portbagaj, şosetă, bici, avion, linie, şoim, personal, har, legătură, scaun, cupru, om de ştiinţă, limuzină} \\
            }
        }
    }
    & \\ \\

    \theutterance \stepcounter{utterance}  
    & & & \multicolumn{4}{p{0.6\linewidth}}{
        \cellcolor[rgb]{0.9,0.9,0.9}{
            \makecell[{{p{\linewidth}}}]{
                \texttt{\tiny{[GM$\langle$P2]}}
                \texttt{GHICIRE: şosetă, personal} \\
            }
        }
    }
     \\ \\

    \theutterance \stepcounter{utterance}  
    & & & \multicolumn{2}{p{0.3\linewidth}}{
        \cellcolor[rgb]{0.9,0.9,0.9}{
            \makecell[{{p{\linewidth}}}]{
                \texttt{\tiny{[GM$|$GM]}}
                ['şosetă', 'personal']
            }
        }
    }
    & & \\ \\

    \theutterance \stepcounter{utterance}  
    & & & \multicolumn{2}{p{0.3\linewidth}}{
        \cellcolor[rgb]{0.9,0.9,0.9}{
            \makecell[{{p{\linewidth}}}]{
                \texttt{\tiny{[GM$|$GM]}}
                {'word': 'şosetă', 'assignment': 'team'}
            }
        }
    }
    & & \\ \\

    \theutterance \stepcounter{utterance}  
    & & & \multicolumn{2}{p{0.3\linewidth}}{
        \cellcolor[rgb]{0.9,0.9,0.9}{
            \makecell[{{p{\linewidth}}}]{
                \texttt{\tiny{[GM$|$GM]}}
                {'word': 'şosetă', 'assignment': 'team'}
            }
        }
    }
    & & \\ \\

    \theutterance \stepcounter{utterance}  
    & & & \multicolumn{2}{p{0.3\linewidth}}{
        \cellcolor[rgb]{0.9,0.9,0.9}{
            \makecell[{{p{\linewidth}}}]{
                \texttt{\tiny{[GM$|$GM]}}
                {'word': 'personal', 'assignment': 'innocent'}
            }
        }
    }
    & & \\ \\

    \theutterance \stepcounter{utterance}  
    & & & \multicolumn{2}{p{0.3\linewidth}}{
        \cellcolor[rgb]{0.9,0.9,0.9}{
            \makecell[{{p{\linewidth}}}]{
                \texttt{\tiny{[GM$|$GM]}}
                \texttt{personal} \\
            }
        }
    }
    & & \\ \\

    \theutterance \stepcounter{utterance}  
    & & & \multicolumn{2}{p{0.3\linewidth}}{
        \cellcolor[rgb]{0.9,0.9,0.9}{
            \makecell[{{p{\linewidth}}}]{
                \texttt{\tiny{[GM$|$GM]}}
                {'word': 'cupru', 'assignment': 'opponent'}
            }
        }
    }
    & & \\ \\

    \theutterance \stepcounter{utterance}  
    & & & \multicolumn{2}{p{0.3\linewidth}}{
        \cellcolor[rgb]{0.9,0.9,0.9}{
            \makecell[{{p{\linewidth}}}]{
                \texttt{\tiny{[GM$|$GM]}}
                {'hidden': {'team': ['har', 'legătură'], 'innocent': ['avion', 'limuzină', 'portbagaj', 'scaun', 'bici', 'şoim'], 'opponent': ['om de ştiinţă', 'linie'], 'assassin': []}, 'revealed': {'team': {'team': ['europa', 'piscină', 'inghetata', 'grevă', 'soldat', 'stâncă', 'şosetă'], 'innocent': ['teren', 'personal'], 'opponent': ['pol'], 'assassin': []}, 'opponent': {'team': [], 'innocent': [], 'opponent': ['laborator', 'maşină de spălat', 'celulă', 'curea', 'cupru'], 'assassin': []}}}
            }
        }
    }
    & & \\ \\

    \theutterance \stepcounter{utterance}  
    & & \multicolumn{4}{p{0.6\linewidth}}{
        \cellcolor[rgb]{0.9,0.9,0.9}{
            \makecell[{{p{\linewidth}}}]{
                \texttt{\tiny{[P1$\langle$GM]}}
                \texttt{Cuvântul şosetă a fost ghicit corect. Cuvântul personal a fost ghicit, dar este un cuvânt nevinovat. Rândul colegului tău de echipă s{-}a încheiat acolo.} \\
\texttt{Acum furnizați un alt indiciu referitor la unele dintre cuvintele echipei dvs. rămase și o listă a cuvintelor țintă aferente. Nu uitați să începeți indiciul cu 'INDICIU: ', puneți un rând nou și începeți lista de cuvinte țintă separate prin virgulă cu 'ȚINTE: '. Observație: unele cuvinte au fost eliminate din liste față de solicitările anterioare.} \\
\\ 
\texttt{Cuvintele echipei dvs. rămase sunt: har, legătură.} \\
\texttt{Cuvintele rămase pentru adversarul dvs. sunt: cupru, om de ştiinţă, linie.} \\
\texttt{Cuvintele de distragere a atenției rămase sunt: avion, limuzină, portbagaj, scaun, bici, şoim.} \\
\texttt{Cuvintele de ucidere rămase sunt: .} \\
            }
        }
    }
    & \\ \\

    \theutterance \stepcounter{utterance}  
    & \multicolumn{4}{p{0.6\linewidth}}{
        \cellcolor[rgb]{0.9,0.9,0.9}{
            \makecell[{{p{\linewidth}}}]{
                \texttt{\tiny{[P1$\rangle$GM]}}
                \texttt{INDICIU: NOD} \\
\texttt{ȚINTE: legătură, har} \\
            }
        }
    }
    & & \\ \\

    \theutterance \stepcounter{utterance}  
    & & & \multicolumn{2}{p{0.3\linewidth}}{
        \cellcolor[rgb]{0.9,0.9,0.9}{
            \makecell[{{p{\linewidth}}}]{
                \texttt{\tiny{[GM$|$GM]}}
                \texttt{nod} \\
            }
        }
    }
    & & \\ \\

    \theutterance \stepcounter{utterance}  
    & & & \multicolumn{2}{p{0.3\linewidth}}{
        \cellcolor[rgb]{0.9,0.9,0.9}{
            \makecell[{{p{\linewidth}}}]{
                \texttt{\tiny{[GM$|$GM]}}
                ['legătură', 'har']
            }
        }
    }
    & & \\ \\

    \theutterance \stepcounter{utterance}  
    & & & \multicolumn{2}{p{0.3\linewidth}}{
        \cellcolor[rgb]{0.9,0.9,0.9}{
            \makecell[{{p{\linewidth}}}]{
                \texttt{\tiny{[GM$|$GM]}}
                {'word': 'legătură', 'assignment': 'team'}
            }
        }
    }
    & & \\ \\

    \theutterance \stepcounter{utterance}  
    & & & \multicolumn{2}{p{0.3\linewidth}}{
        \cellcolor[rgb]{0.9,0.9,0.9}{
            \makecell[{{p{\linewidth}}}]{
                \texttt{\tiny{[GM$|$GM]}}
                {'word': 'har', 'assignment': 'team'}
            }
        }
    }
    & & \\ \\

    \theutterance \stepcounter{utterance}  
    & & \multicolumn{4}{p{0.6\linewidth}}{
        \cellcolor[rgb]{0.9,0.9,0.9}{
            \makecell[{{p{\linewidth}}}]{
                \texttt{\tiny{[GM$\rangle$P2]}}
                \texttt{Cuvântul şosetă a fost ghicit corect. Cuvântul personal a fost ghicit, dar este un cuvânt nevinovat. Rândul tău s{-}a încheiat acolo.} \\
\texttt{Acum furnizați o altă listă separată prin virgulă de cel puțin 1 și până la 2 cuvinte din următoarea listă de cuvinte care se referă cel mai bine sau sunt cel mai strâns asociate cu cuvântul 'nod'. Nu uitați să vă începeți răspunsul cu GHICIRE:. Notă: unele cuvinte au fost eliminate din listă față de solicitările anterioare.} \\
\\ 
\texttt{portbagaj, bici, avion, linie, şoim, har, legătură, scaun, om de ştiinţă, limuzină} \\
            }
        }
    }
    & \\ \\

    \theutterance \stepcounter{utterance}  
    & & & \multicolumn{4}{p{0.6\linewidth}}{
        \cellcolor[rgb]{0.9,0.9,0.9}{
            \makecell[{{p{\linewidth}}}]{
                \texttt{\tiny{[GM$\langle$P2]}}
                \texttt{GHICIRE: legătură, linie} \\
            }
        }
    }
     \\ \\

    \theutterance \stepcounter{utterance}  
    & & & \multicolumn{2}{p{0.3\linewidth}}{
        \cellcolor[rgb]{0.9,0.9,0.9}{
            \makecell[{{p{\linewidth}}}]{
                \texttt{\tiny{[GM$|$GM]}}
                ['legătură', 'linie']
            }
        }
    }
    & & \\ \\

    \theutterance \stepcounter{utterance}  
    & & & \multicolumn{2}{p{0.3\linewidth}}{
        \cellcolor[rgb]{0.9,0.9,0.9}{
            \makecell[{{p{\linewidth}}}]{
                \texttt{\tiny{[GM$|$GM]}}
                {'word': 'legătură', 'assignment': 'team'}
            }
        }
    }
    & & \\ \\

    \theutterance \stepcounter{utterance}  
    & & & \multicolumn{2}{p{0.3\linewidth}}{
        \cellcolor[rgb]{0.9,0.9,0.9}{
            \makecell[{{p{\linewidth}}}]{
                \texttt{\tiny{[GM$|$GM]}}
                {'word': 'legătură', 'assignment': 'team'}
            }
        }
    }
    & & \\ \\

    \theutterance \stepcounter{utterance}  
    & & & \multicolumn{2}{p{0.3\linewidth}}{
        \cellcolor[rgb]{0.9,0.9,0.9}{
            \makecell[{{p{\linewidth}}}]{
                \texttt{\tiny{[GM$|$GM]}}
                {'word': 'linie', 'assignment': 'opponent'}
            }
        }
    }
    & & \\ \\

    \theutterance \stepcounter{utterance}  
    & & & \multicolumn{2}{p{0.3\linewidth}}{
        \cellcolor[rgb]{0.9,0.9,0.9}{
            \makecell[{{p{\linewidth}}}]{
                \texttt{\tiny{[GM$|$GM]}}
                \texttt{linie} \\
            }
        }
    }
    & & \\ \\

    \theutterance \stepcounter{utterance}  
    & & & \multicolumn{2}{p{0.3\linewidth}}{
        \cellcolor[rgb]{0.9,0.9,0.9}{
            \makecell[{{p{\linewidth}}}]{
                \texttt{\tiny{[GM$|$GM]}}
                {'word': 'om de ştiinţă', 'assignment': 'opponent'}
            }
        }
    }
    & & \\ \\

    \theutterance \stepcounter{utterance}  
    & & & \multicolumn{2}{p{0.3\linewidth}}{
        \cellcolor[rgb]{0.9,0.9,0.9}{
            \makecell[{{p{\linewidth}}}]{
                \texttt{\tiny{[GM$|$GM]}}
                {'hidden': {'team': ['har'], 'innocent': ['avion', 'limuzină', 'portbagaj', 'scaun', 'bici', 'şoim'], 'opponent': [], 'assassin': []}, 'revealed': {'team': {'team': ['europa', 'piscină', 'inghetata', 'grevă', 'soldat', 'stâncă', 'şosetă', 'legătură'], 'innocent': ['teren', 'personal'], 'opponent': ['pol', 'linie'], 'assassin': []}, 'opponent': {'team': [], 'innocent': [], 'opponent': ['laborator', 'maşină de spălat', 'celulă', 'curea', 'cupru', 'om de ştiinţă'], 'assassin': []}}}
            }
        }
    }
    & & \\ \\

    \theutterance \stepcounter{utterance}  
    & & \multicolumn{4}{p{0.6\linewidth}}{
        \cellcolor[rgb]{0.9,0.9,0.9}{
            \makecell[{{p{\linewidth}}}]{
                \texttt{\tiny{[P1$\langle$GM]}}
                \texttt{Cuvântul legătură a fost ghicit corect. Cuvântul linie a fost ghicit, dar este un cuvânt adversar. Rândul colegului tău de echipă s{-}a încheiat acolo.} \\
\texttt{Acum furnizați un alt indiciu referitor la unele dintre cuvintele echipei dvs. rămase și o listă a cuvintelor țintă aferente. Nu uitați să începeți indiciul cu 'INDICIU: ', puneți un rând nou și începeți lista de cuvinte țintă separate prin virgulă cu 'ȚINTE: '. Observație: unele cuvinte au fost eliminate din liste față de solicitările anterioare.} \\
\\ 
\texttt{Cuvintele echipei dvs. rămase sunt: har.} \\
\texttt{Cuvintele rămase pentru adversarul dvs. sunt: om de ştiinţă.} \\
\texttt{Cuvintele de distragere a atenției rămase sunt: avion, limuzină, portbagaj, scaun, bici, şoim.} \\
\texttt{Cuvintele de ucidere rămase sunt: .} \\
            }
        }
    }
    & \\ \\

    \theutterance \stepcounter{utterance}  
    & \multicolumn{4}{p{0.6\linewidth}}{
        \cellcolor[rgb]{0.9,0.9,0.9}{
            \makecell[{{p{\linewidth}}}]{
                \texttt{\tiny{[P1$\rangle$GM]}}
                \texttt{INDICIU: SORTE} \\
\texttt{ȚINTE: har} \\
            }
        }
    }
    & & \\ \\

    \theutterance \stepcounter{utterance}  
    & & & \multicolumn{2}{p{0.3\linewidth}}{
        \cellcolor[rgb]{0.9,0.9,0.9}{
            \makecell[{{p{\linewidth}}}]{
                \texttt{\tiny{[GM$|$GM]}}
                \texttt{sorte} \\
            }
        }
    }
    & & \\ \\

    \theutterance \stepcounter{utterance}  
    & & & \multicolumn{2}{p{0.3\linewidth}}{
        \cellcolor[rgb]{0.9,0.9,0.9}{
            \makecell[{{p{\linewidth}}}]{
                \texttt{\tiny{[GM$|$GM]}}
                ['har']
            }
        }
    }
    & & \\ \\

    \theutterance \stepcounter{utterance}  
    & & & \multicolumn{2}{p{0.3\linewidth}}{
        \cellcolor[rgb]{0.9,0.9,0.9}{
            \makecell[{{p{\linewidth}}}]{
                \texttt{\tiny{[GM$|$GM]}}
                {'word': 'har', 'assignment': 'team'}
            }
        }
    }
    & & \\ \\

    \theutterance \stepcounter{utterance}  
    & & & \multicolumn{2}{p{0.3\linewidth}}{
        \cellcolor[rgb]{0.9,0.9,0.9}{
            \makecell[{{p{\linewidth}}}]{
                \texttt{\tiny{[GM$|$GM]}}
                \texttt{opponent has won} \\
            }
        }
    }
    & & \\ \\

\end{supertabular}
}

\end{document}
